%%%%%%%%%%%%%%%%%%%%%%%%%
\spellentry{Ghost Sound}
%%%%%%%%%%%%%%%%%%%%%%%%%

\tagline{Illusion (Figment)}
\begin{description*}
\item[Level:] Brd 0, Sor/Wiz 0
\item[Casting Time:] 1 standard action
\item[Range:] Close (25 ft. + 5 ft./2 levels)
\item[Effect:] Illusory sounds
\item[Duration:] 1 round/level (D)
\item[Saving Throw:] Will disbelief (if interacted with)
\item[Spell Resistance:] No
\item[Components:] V, S, M
\item[Arcane Material Component:] A bit of wool or a small lump of wax.
\end{description*}

\textit{Ghost sound} allows you to create a volume of sound that rises, recedes,
approaches, or remains at a fixed place. You choose what type of sound \textit{ghost
sound} creates when casting it and cannot thereafter change the sound's basic character.
The volume of sound created depends on your level. You can produce as much noise
as four normal humans per caster level (maximum twenty humans). Thus, talking,
singing, shouting, walking, marching, or running sounds can be created. The noise
a \textit{ghost sound} spell produces can be virtually any type of sound within
the volume limit. A horde of rats running and squeaking is about the same volume
as eight humans running and shouting. A roaring lion is equal to the noise from
sixteen humans, while a roaring dire tiger is equal to the noise from twenty humans.
\textit{Ghost sound} can enhance the effectiveness of a \textit{silent image} spell.
\textit{Ghost sound} can be made permanent with a \textit{permanency} spell.